\documentclass[12pt,a4paper]{article}

\usepackage[T2A]{fontenc}
\usepackage[utf8]{inputenc}
\usepackage[english,russian]{babel}
\usepackage{cmap}
\usepackage{amsmath,amssymb,mathtools}
\usepackage{geometry}
\usepackage{microtype}
\usepackage{enumitem}
\usepackage{titlesec}
\usepackage{hyperref}
\usepackage{bookmark}
\usepackage{tikz}
\usetikzlibrary{arrows.meta,calc}

\geometry{left=20mm,right=20mm,top=22mm,bottom=22mm}

\setlength{\parindent}{0pt}
\setlength{\parskip}{5pt}
\setcounter{tocdepth}{2}

\titleformat{\section}
  {\Large\bfseries}{\thesection}{0.7em}{}
\titleformat{\subsection}
  {\large\bfseries}{\thesubsection}{0.7em}{}

\hypersetup{
  unicode=true,
  colorlinks=true,
  linkcolor=black,
  urlcolor=black,
  pdftitle={ТФКП. Лекция 1}
}

\newcommand{\CC}{\mathbb C}
\newcommand{\RR}{\mathbb R}
\newcommand{\OO}{\mathcal O}

\begin{document}

% ------------------------------------------------------------
% Титульный лист
% ------------------------------------------------------------
\begin{titlepage}
\begin{center}

\vspace*{1.5cm}

{\LARGE\bfseries Теория функций комплексного переменного}

\vspace{1.2cm}

{\large Механико-математический факультет МГУ}

\vspace{2cm}

{\Large Конспект лекций}

\vspace{1.5cm}

{\large Лектор: Андрей Викторович Домрин}

\vspace{0.7cm}

{\large 5 семестр}

\vfill

{\Large\bfseries Лекция 1}

\vspace{1.5cm}

\end{center}
\end{titlepage}

% ------------------------------------------------------------
% Содержание
% ------------------------------------------------------------
\tableofcontents
\newpage

% ------------------------------------------------------------
\section{Лекция 1}
% ------------------------------------------------------------

\subsection{Комплексные числа}

Напомним основные обозначения. Комплексное число
\(z=x+iy\in\CC\) отождествляется с точкой \((x,y)\in\RR^2\). При этом
\[
\operatorname{Re}z=x,\qquad \operatorname{Im}z=y.
\]

Модуль комплексного числа равен
\[
|z|=\sqrt{x^2+y^2}.
\]
Если \(z\neq0\), то
\[
z=|z|(\cos\varphi+i\sin\varphi),
\]
где \(\varphi=\arg z\). Аргумент определён с точностью до \(2\pi\).

\begin{center}
\begin{tikzpicture}[>=Latex,scale=0.9]
  \draw[->] (-0.3,0)--(4.2,0) node[right] {$\operatorname{Re}z$};
  \draw[->] (0,-0.3)--(0,3.2) node[above] {$\operatorname{Im}z$};
  \coordinate (O) at (0,0);
  \coordinate (Z) at (3,2.1);
  \draw[thick,->] (O)--(Z) node[midway,above left] {$|z|$};
  \draw (0.9,0) arc[start angle=0,end angle=35,radius=0.9];
  \node at (1.18,0.34) {$\arg z$};
  \fill (Z) circle (1.4pt);
  \node[above right] at (Z) {$z=x+iy$};
\end{tikzpicture}
\end{center}

Пусть \(z_1,z_2\in\CC\). Тогда
\[
|z_1z_2|=|z_1|\,|z_2|,
\]
а аргументы при умножении складываются с точностью до \(2\pi\):
\[
\arg(z_1z_2)=\arg z_1+\arg z_2 \pmod{2\pi}.
\]

\textbf{Пример.} Найти \(A,B\in\RR\), если
\[
(1+i)^{2026}=A+iB.
\]

\textbf{Решение.} Так как
\[
1+i=\sqrt2\left(\cos\frac{\pi}{4}+i\sin\frac{\pi}{4}\right),
\]
то
\[
(1+i)^{2026}
=2^{1013}\left(\cos\frac{2026\pi}{4}+i\sin\frac{2026\pi}{4}\right).
\]
Поскольку
\[
\frac{2026\pi}{4}=506\pi+\frac{\pi}{2},
\]
получаем
\[
(1+i)^{2026}=2^{1013}i.
\]
Следовательно,
\[
A=0,\qquad B=2^{1013}.
\]

% ------------------------------------------------------------
\subsection{Непрерывность функции комплексного переменного}

Пусть \(D\subset\CC\) --- открытое множество и \(f:D\to\CC\).

\textbf{Определение.} Функция \(f\) называется непрерывной в точке
\(z_0\in D\), если
\[
\forall\varepsilon>0\ \exists\delta>0:\quad
|z-z_0|<\delta\ \Longrightarrow\ |f(z)-f(z_0)|<\varepsilon.
\]

Эквивалентно: для любой последовательности \(z_n\to z_0\)
\[
f(z_n)\to f(z_0).
\]

\begin{center}
\begin{tikzpicture}[>=Latex,scale=0.82]
  \draw[->] (-3.7,0)--(-0.8,0);
  \draw[->] (-2.25,-1.45)--(-2.25,1.65);
  \draw[thick] (-2.25,0.05) ellipse (1.0 and 0.7);
  \fill (-2.25,0.05) circle (1.4pt) node[below] {$z_0$};
  \draw (-2.25,0.05) circle (0.38);
  \node at (-2.25,1.08) {$D$};

  \draw[->,thick] (-0.45,0.2)--(0.85,0.2) node[midway,above] {$w=f(z)$};

  \draw[->] (1.2,0)--(4.2,0);
  \draw[->] (2.7,-1.45)--(2.7,1.65);
  \draw[thick] (2.7,0.05) ellipse (1.0 and 0.7);
  \fill (2.7,0.05) circle (1.4pt) node[below] {$f(z_0)$};
  \draw (2.7,0.05) circle (0.38);
\end{tikzpicture}
\end{center}

\textbf{Примеры.}
\begin{enumerate}[label=\arabic*)]
  \item \(f(z)=z\) непрерывна на \(\CC\).
  \item \(f(z)=z^2\) непрерывна на \(\CC\).
  \item \(f(z)=\overline z\) непрерывна на \(\CC\).
  \item \(f(z)=1/z\) определена и непрерывна на \(\CC\setminus\{0\}\).
  \item Для функции
  \[
  f(z)=\ln|z|+i\arg z
  \]
  необходимо выбрать ветвь аргумента.
\end{enumerate}

Если выбрать аргумент так, что \(-\pi<\arg z\leq\pi\), то на отрицательной
вещественной полуоси возникает разрыв. Действительно, при подходе к точке
\(-1\) сверху аргументы стремятся к \(\pi\), а при подходе снизу --- к
\(-\pi\). Поэтому для однозначного непрерывного выбора аргумента требуется
разрез плоскости.

% ------------------------------------------------------------
\subsection{Дифференцируемость и условия Коши--Римана}

Пусть \(D\subset\CC\) --- открытое множество, \(z_0\in D\),
\(f:D\to\CC\).

\textbf{Определение.} Функция \(f\) называется комплексно дифференцируемой
в точке \(z_0\), если существует предел
\[
f'(z_0)=\lim\limits_{z\to z_0}
\frac{f(z)-f(z_0)}{z-z_0}\in\CC.
\]

Эквивалентно, существует \(A\in\CC\) такое, что
\[
f(z_0+h)=f(z_0)+Ah+o(|h|),\qquad h\to0.
\]
При этом \(A=f'(z_0)\).

Теперь будем рассматривать \(f\) как отображение \(\RR^2\to\RR^2\). Пусть
\[
f(z)=u(x,y)+iv(x,y),\qquad z=x+iy.
\]

\textbf{Определение.} Функция \(f\) называется вещественно дифференцируемой
в точке \(z_0\), если существует линейное отображение
\(d_{z_0}f:\RR^2\to\RR^2\), для которого
\[
f(z_0+h)-f(z_0)=d_{z_0}f(h)+o(|h|),\qquad h\to0.
\]
Если \(f\) вещественно дифференцируема, то
\[
d_{z_0}f=
\begin{pmatrix}
 u_x(z_0) & u_y(z_0)\\
 v_x(z_0) & v_y(z_0)
\end{pmatrix}.
\]

\textbf{Теорема (условия Коши--Римана).}
Функция \(f=u+iv\) комплексно дифференцируема в точке \(z_0\) тогда и
только тогда, когда она вещественно дифференцируема в точке \(z_0\) и
выполнены условия
\[
u_x(z_0)=v_y(z_0),\qquad
u_y(z_0)=-v_x(z_0).
\]

\textbf{Доказательство.}
Пусть \(f\) комплексно дифференцируема в \(z_0\). Тогда
\[
f(z_0+h)=f(z_0)+Ah+o(|h|),
\]
где \(A=a+ib\), \(h=\xi+i\eta\). Имеем
\[
Ah=(a+ib)(\xi+i\eta)
=(a\xi-b\eta)+i(b\xi+a\eta).
\]
Значит, умножению на \(A\) как вещественному линейному отображению
соответствует матрица
\[
\begin{pmatrix}
 a & -b\\
 b & a
\end{pmatrix}.
\]
Следовательно,
\[
\begin{pmatrix}
 u_x(z_0) & u_y(z_0)\\
 v_x(z_0) & v_y(z_0)
\end{pmatrix}
=
\begin{pmatrix}
 a & -b\\
 b & a
\end{pmatrix},
\]
откуда и следуют условия Коши--Римана.

Обратно, если \(f\) вещественно дифференцируема и выполнены условия
Коши--Римана, то матрица \(d_{z_0}f\) имеет вид
\[
\begin{pmatrix}
 a & -b\\
 b & a
\end{pmatrix},
\]
то есть соответствует умножению на комплексное число \(A=a+ib\). Поэтому
\[
f(z_0+h)=f(z_0)+Ah+o(|h|),
\]
и функция комплексно дифференцируема в \(z_0\).

\textbf{Следствие.} Если \(f=u+iv\) комплексно дифференцируема в \(z_0\), то
\[
f'(z_0)
=u_x(z_0)+iv_x(z_0)
=v_y(z_0)-iu_y(z_0).
\]

% ------------------------------------------------------------
\subsection{Голоморфные функции}

\textbf{Определение.} Функция \(f:D\to\CC\) называется голоморфной в точке
\(z_0\in D\), если она комплексно дифференцируема в некоторой окрестности
этой точки. Запись:
\[
f\in\OO(z_0).
\]

Функция называется голоморфной на открытом множестве \(D\), если она
комплексно дифференцируема в каждой точке \(D\). Запись:
\[
f\in\OO(D).
\]

Для произвольного множества \(M\subset\CC\) запись \(f\in\OO(M)\) означает,
что существует открытое множество \(D\supset M\), на котором \(f\) голоморфна.

\textbf{Свойства.}
\begin{enumerate}[label=\arabic*)]
  \item Если \(f,g\in\OO(D)\), то \(f\pm g\) и \(fg\) голоморфны на \(D\), причём
  \[
  (fg)'=f'g+fg'.
  \]

  \item Если \(g(z)\neq0\), то частное \(f/g\) голоморфно там, где знаменатель
  не обращается в нуль.

  \item Если \(f\in\OO(D_1)\), \(g\in\OO(D_2)\) и \(D_2\supset f(D_1)\), то
  \(g\circ f\in\OO(D_1)\), причём
  \[
  (g\circ f)'(z)=g'(f(z))f'(z).
  \]
\end{enumerate}

% ------------------------------------------------------------
\subsection{Обратная функция}

\textbf{Теорема.} Пусть \(P_1,P_2\subset\CC\) --- открытые множества,
\(f\in\OO(P_1)\) гомеоморфно отображает \(P_1\) на \(P_2\), и
\[
f'(z)\neq0\qquad \forall z\in P_1.
\]
Тогда обратное отображение \(g=f^{-1}:P_2\to P_1\) голоморфно, причём
\[
g'(w)=\frac{1}{f'(g(w))},\qquad w\in P_2.
\]

\textbf{Доказательство.}
Пусть \(w=f(z)\), \(w_1=f(z_1)\). Тогда
\[
\frac{g(w_1)-g(w)}{w_1-w}
=
\frac{z_1-z}{f(z_1)-f(z)}.
\]
Так как \(g\) непрерывна, из \(w_1\to w\) следует \(z_1\to z\). Поэтому
\[
\frac{g(w_1)-g(w)}{w_1-w}
\longrightarrow
\frac{1}{f'(z)}
=
\frac{1}{f'(g(w))}.
\]

\textbf{Замечание.} На лекции отмечено, что достаточно требовать, чтобы
голоморфная функция биективно отображала \(P_1\) на \(P_2\): необходимые
свойства обратного отображения в этой ситуации следуют автоматически.

% ------------------------------------------------------------
\subsection{Экспонента и комплексный логарифм}

Рассмотрим комплексную экспоненту
\[
w=e^z.
\]
Её можно задавать степенным рядом
\[
e^z=\sum\limits_{n=0}^{\infty}\frac{z^n}{n!}.
\]
Для \(z=x+iy\) используется представление
\[
e^{x+iy}=e^x(\cos y+i\sin y).
\]
В частности,
\[
e^{iy}=\cos y+i\sin y.
\]

Для любых \(z_1,z_2\in\CC\)
\[
e^{z_1+z_2}=e^{z_1}e^{z_2}.
\]

Покажем, что \(f(z)=e^z\) голоморфна. Пишем
\[
u(x,y)=e^x\cos y,\qquad v(x,y)=e^x\sin y.
\]
Тогда
\[
u_x=e^x\cos y=v_y,
\qquad
u_y=-e^x\sin y=-v_x.
\]
Условия Коши--Римана выполнены на всей \(\CC\). Поэтому
\[
e^z\in\OO(\CC),
\qquad
(e^z)'=e^z.
\]

Вернёмся к нескольким примерам:
\[
f_1(z)=z,\qquad f_1'(z)=1,
\]
\[
f_2(z)=z^2,\qquad f_2'(z)=2z.
\]
Для \(f_3(z)=\overline z\) имеем \(u=x\), \(v=-y\), поэтому
\[
d_zf_3=
\begin{pmatrix}
1&0\\
0&-1
\end{pmatrix}.
\]
Эта матрица не имеет вида
\(\begin{pmatrix}a&-b\\ b&a\end{pmatrix}\), поэтому функция
\(\overline z\) нигде не комплексно дифференцируема.

Кроме того,
\[
f_4(z)=\frac1z\in\OO(\CC\setminus\{0\}),
\qquad
f_4'(z)=-\frac1{z^2}.
\]

Теперь рассмотрим область
\[
D=\CC\setminus(-\infty,0].
\]
На ней выбирается ветвь аргумента
\[
-\pi<\arg z<\pi.
\]

\textbf{Определение.} Соответствующая ветвь комплексного логарифма задаётся
формулой
\[
\operatorname{Log}z=\ln|z|+i\arg z.
\]

Проверим:
\[
e^{\operatorname{Log}z}
=e^{\ln|z|+i\arg z}
=|z|\bigl(\cos(\arg z)+i\sin(\arg z)\bigr)
=z.
\]
Таким образом, эта ветвь логарифма является обратной к экспоненте после
ограничения экспоненты на горизонтальную полосу
\[
-\pi<\operatorname{Im}w<\pi.
\]

\begin{center}
\begin{tikzpicture}[>=Latex,scale=0.78]
  \draw[->] (-4,0)--(-0.6,0) node[right] {$x$};
  \draw[->] (-2.3,-1.7)--(-2.3,1.7) node[above] {$y$};
  \draw[line width=1.2pt] (-4,0)--(-2.3,0);
  \node at (-2.3,-2.05) {$D=\CC\setminus(-\infty,0]$};

  \draw[->,thick] (-0.3,0)--(1.0,0) node[midway,above] {$w=\operatorname{Log}z$};

  \draw[->] (1.4,0)--(5.1,0) node[right] {$u$};
  \draw[->] (3.1,-1.9)--(3.1,1.9) node[above] {$v$};
  \draw[dashed] (1.6,1.25)--(4.9,1.25);
  \draw[dashed] (1.6,-1.25)--(4.9,-1.25);
  \node[right] at (4.9,1.25) {$\pi$};
  \node[right] at (4.9,-1.25) {$-\pi$};
\end{tikzpicture}
\end{center}

По формуле для производной обратной функции
\[
(\operatorname{Log}z)'
=\frac{1}{e^{\operatorname{Log}z}}
=\frac1z.
\]

\end{document}
